\documentclass[11pt,a4paper]{article}

\usepackage[T1]{fontenc}
\usepackage[utf8]{inputenc}
\usepackage{lmodern}
\usepackage{xcolor}
\usepackage{amsmath}
\usepackage{amssymb}
\usepackage{geometry}
\geometry{margin=1in}

\definecolor{linkblue}{RGB}{0,86,179}
\usepackage[
    colorlinks=true,
    linkcolor=linkblue,
    citecolor=linkblue,
    urlcolor=linkblue,
    pdftitle={Geotemporal Hydrodynamics: The Unified Fluid Theory},
    pdfauthor={T. Abram},
    pdfsubject={Physics, Cosmology},
    pdfkeywords={Fluid Theory, Gravity, Cosmology},
    bookmarks=true,
    pdfdisplaydoctitle=true,
]{hyperref}

\title{Geotemporal Hydrodynamics\\ \Large The Unified Fluid Theory}
\author{T. Abram}
\date{March 2026}

\begin{document}

\maketitle

\begin{abstract}
Geotemporal Hydrodynamics (GTH) proposes that spacetime is a macroscopic quantum superfluid
rather than an empty geometric manifold. This framework models the vacuum as a 5D
viscoelastic substrate with order parameter $\Psi = \rho e^{i\theta}$. This new report highlights four
key phenomenological pillars of the theory based on recent visualizations.
\end{abstract}

\section{Fluid Wake and Flow}
Gravity emerges as the long-wavelength acoustic response of the substrate. The effective flow velocity incorporates both a Newtonian gravitational term and a macroscopic wake term parameterized by the Reynolds number ($Re$) and fluid density ($\rho_0$).
\begin{equation}
v = \sqrt{\frac{GM}{r} + \beta r}
\end{equation}
where $\beta \propto Re \cdot \rho_0$.

\section{GeoKnots (Topology)}
Particles are not point objects but stable topological defects (GeoKnots) in the 5D viscoelastic substrate. The fundamental particle is modeled as a Trefoil Knot ($3_1$). Rest mass is defined as the localized kinetic bending energy of the substrate:
\begin{equation}
m_e \propto \frac{\sqrt{\rho_0}}{m_\psi^2}
\end{equation}
Adjusting the scalar mass ($m_\psi$) changes the knot's thickness, while $\rho_0$ alters its overall scale.

\section{Viscoelastic Relaxation}
The substrate exhibits finite relaxation times ($\tau$) when deformed, leading to viscoelastic damping. The fluid's response to rapid perturbations is governed by Israel-Stewart hydrodynamics, ensuring causality and providing a mechanism for ringdown dispersion/damping consistent with gravitational wave observations.

\section{Rotation Curves}
At galactic scales, the inverse energy cascade of the superfluid vacuum leads to anomalous rotation curves. While Newtonian predictions fall off as $v \propto 1/\sqrt{r}$, the GTH framework provides a non-Keplerian flat profile at large radii due to the $\beta r$ term, naturally accounting for dark matter phenomenology.

\end{document}
